%% The Merge witness against per-side naming of conversion leaves
%% (2026-10-26).  Examples §1b, run by strong-rep-nu's evaluator; the
%% checked version is notes/SideProbe.agda §5 (`flawsAt 5`, `flawsAt 6`,
%% `rawAt 6`).  Notation follows paper/main.tex: [δ] M ⟨c⟩ is a boundary,
%% +X^α / -X^α are scope entries (bind / unbind, acting left to right),
%% and -X / +X are the conversions seal / unseal.
%%
%% Build: pdflatex MergeWitness.tex
\documentclass[11pt]{article}
\usepackage[margin=0.8in]{geometry}
\usepackage{amsmath,amssymb}
\usepackage{stmaryrd}
\allowdisplaybreaks

\newcommand{\key}[1]{\mathtt{#1}}
\newcommand{\id}[1]{\key{id}(#1)}
\newcommand{\seal}[1]{-#1}
\newcommand{\unseal}[1]{+#1}
\newcommand{\becomes}{\longrightarrow}
\newcommand{\of}{{:}}
\newcommand{\tu}{{\to}}
\newcommand{\dotspace}{.\,}
\newcommand{\lam}[1]{\lambda #1 \dotspace}
\newcommand{\Lam}[1]{\Lambda #1 \dotspace}
\newcommand{\all}[1]{\forall #1 \dotspace}
\newcommand{\app}{\;}
\newcommand{\Bool}{\key{Bool}}
\newcommand{\True}{\key{true}}
\newcommand{\False}{\key{false}}
\newcommand{\by}{{:=}}
\newcommand{\nuinst}[4]{\nu #1\by #2 \dotspace (#3 \app #1) \langle #4 \rangle}
\newcommand{\boundry}[3]{\left[#2\right] #1 \langle #3 \rangle}
\newcommand{\cmp}{\mathbin{\fatsemi}}
\newcommand{\bnd}[1]{{+#1}}
\newcommand{\ubnd}[1]{{-#1}}
\newcommand{\rul}[1]{\textsc{#1}}

\begin{document}

\begin{center}
  {\Large A \textsc{Merge} that loses a name}\\[4pt]
  strong-rep-nu, Examples \S1b; checked in
  \texttt{notes/SideProbe.agda} \S5
\end{center}

\section*{The question}

In the \emph{per-side} design, a seal or unseal names its type
variable on the side where the scope gives it a name, instead of in the
conversion (union) context $\delta^{+}(\Delta)$.  For a conversion at
polarity $p$ in a boundary $\boundry{M}{\delta}{c}$ with exterior
$\Delta$ and interior $\delta(\Delta)$:
\begin{itemize}
\item $\seal{X}$ at $+$ is named on the exterior, at $-$ on the interior;
\item $\unseal{X}$ at $+$ is named on the interior, at $-$ on the exterior.
\end{itemize}
In a chain, each later link's source is the previous link's target, so
every type in the middle of a chain must be named on the chain's side.
The witness below reaches a boundary whose seal chain passes through a
type variable that \emph{neither} side names.

\section*{The program}

The closed System F source program is
\[
  \bigl((\Lam{X} \lam{f \of \all{W} W \tu \Bool} f\,[X])\,[\Bool]
     \app (\Lam{Z} \lam{z \of Z} \True)\bigr) \app \False
\]
Each type application compiles to a $\nu$ whose conversion reveals the
argument, so the program compiles to state~0 below, writing
\[
  U \;=\; \lam{f \of \all{W} W \tu \Bool}
    \nuinst{Y}{X}{f}{\seal{Y} \tu \id{\Bool}}
  \qquad
  P \;=\; \Lam{Z} \lam{z \of Z} \True
\]
for the body of the type abstraction over $X$ and for the polymorphic
argument.

\section*{The run}

The store records representations: two cells are allocated, $\alpha$ for
$X$ and then $\beta$ for $Y$.  Every step but the first fires under a
frame ($\xi$); the context at the hole is noted.

\begin{align*}
&\text{0.}\quad
  \bigl(\nuinst{X}{\Bool}{(\Lam{X} U)}
     {(\all{W} \id{W} \tu \id{\Bool}) \tu (\seal{X} \tu \id{\Bool})}
   \app P\bigr) \app \False
\\[2pt]
&\quad\becomes\; \text{\rul{TyBeta}, allocating } \alpha \by \Bool
\\[2pt]
&\text{1.}\quad
  \bigl(\boundry{U}{\bnd{X}^\alpha}
     {(\all{W} \id{W} \tu \id{\Bool}) \tu (\seal{X} \tu \id{\Bool})}
   \app P\bigr) \app \False
\\[2pt]
&\quad\becomes\; \text{\rul{Wrap}: } P \text{ enters } X\text{'s scope
  under the dual entry } \ubnd{X}^\alpha
\\[2pt]
&\text{2.}\quad
  \boundry{\bigl(U \app
     \boundry{P}{\ubnd{X}^\alpha}{\all{W} \id{W} \tu \id{\Bool}}\bigr)}
   {\bnd{X}^\alpha}{\seal{X} \tu \id{\Bool}} \app \False
\\[2pt]
&\quad\becomes\; \text{\rul{Beta}, at } X \by \alpha
\\[2pt]
&\text{3.}\quad
  \boundry{\Bigl(\nuinst{Y}{X}
     {\boundry{P}{\ubnd{X}^\alpha}{\all{W} \id{W} \tu \id{\Bool}}}
     {\seal{Y} \tu \id{\Bool}}\Bigr)}
   {\bnd{X}^\alpha}{\seal{X} \tu \id{\Bool}} \app \False
\\[2pt]
&\quad\becomes\; \text{\rul{TyWrap}, at } X \by \alpha,
  \text{ allocating } \beta \by \alpha
  \quad (\text{the payload of } Y \text{ is } X\text{'s cell})
\\[2pt]
&\text{4.}\quad
  \boundry{\Bigl(\boundry
     {\bigl(\boundry{(\lam{z \of Y} \True)}{\ubnd{X}^\alpha}
        {\id{Y} \tu \id{\Bool}}\bigr)}
     {\bnd{Y}^\beta}{\seal{Y} \tu \id{\Bool}}\Bigr)}
   {\bnd{X}^\alpha}{\seal{X} \tu \id{\Bool}} \app \False
\\[2pt]
&\quad\becomes\; \text{\rul{Merge}, at } X \by \alpha:\;
  (\id{Y} \tu \id{\Bool}) \cmp (\seal{Y} \tu \id{\Bool})
  = \seal{Y} \tu \id{\Bool}
\\[2pt]
&\text{5.}\quad
  \boundry{\bigl(\boundry{(\lam{z \of Y} \True)}
     {\bnd{Y}^\beta, \ubnd{X}^\alpha}{\seal{Y} \tu \id{\Bool}}\bigr)}
   {\bnd{X}^\alpha}{\seal{X} \tu \id{\Bool}} \app \False
\\[2pt]
&\quad\becomes\; \text{\rul{Merge}:}\;
  (\seal{Y} \tu \id{\Bool}) \cmp (\seal{X} \tu \id{\Bool})
  = (\seal{X} \cmp \seal{Y}) \tu \id{\Bool}
  = (\seal{X} ; \seal{Y}) \tu \id{\Bool}
\\[2pt]
&\text{6.}\quad
  \boundry{(\lam{z \of Y} \True)}
     {\bnd{X}^\alpha, \bnd{Y}^\beta, \ubnd{X}^\alpha}
     {(\seal{X} ; \seal{Y}) \tu \id{\Bool}}
   \app \False
\end{align*}
The store is $\alpha \by \Bool,\ \beta \by \alpha$ from state~4 on.
(The run goes on by \rul{Wrap}, \rul{Beta} and \rul{Id} to $\True$.)

\section*{State 5 passes}

In state~5 the outer boundary has exterior $\emptyset$ and interior
$X \by \alpha$; its domain $\seal{X}$ sits at polarity $-$, so it is
named on the interior ($X \by \alpha$), and its source $\Bool$ is
read on the exterior.  The inner boundary has exterior $X \by \alpha$
and scope $\bnd{Y}^\beta, \ubnd{X}^\alpha$, so its interior is
$Y \by \beta$; its domain $\seal{Y}$ at $-$ is named on the interior
($Y \by \beta$), and its source, $\beta$'s payload $\alpha$, is read on
the exterior as $X$.  Every leaf has its name on the right side.

\section*{State 6 fails}

\begin{center}
\renewcommand{\arraystretch}{1.3}
\begin{tabular}{ll}
  scope $\delta$ & $\bnd{X}^\alpha,\ \bnd{Y}^\beta,\ \ubnd{X}^\alpha$
    \quad (acting left to right) \\
  exterior $\Delta$ & $\emptyset$ \\
  interior $\delta(\Delta)$ & $Y \by \beta$ \\
  union $\delta^{+}(\Delta)$ & $X \by \alpha,\ Y \by \beta$ \\
  store & $\alpha \by \Bool,\ \beta \by \alpha$
\end{tabular}
\end{center}

The domain is the seal chain $\seal{X} ; \seal{Y}$ at polarity $-$:
\[
  \Bool \;\xrightarrow{\;\seal{X}\;}\; X
        \;\xrightarrow{\;\seal{Y}\;}\; Y
\]
Its seals belong on the interior, which names only $Y$.  The middle type
$X$ (that is, $\alpha$) has a name only in the union: the merged scope
binds $X$ and then unbinds it.  So $\seal{X}$ has no per-side
spelling, and neither does $\seal{Y}$'s source.  The two
\textsc{Merge}s fused three frames, and the one frame that named $X$
on its own side ($\bnd{X}^\alpha$, state~5's outer boundary) is no longer
a side of anything.

\section*{Where it comes from}

Five of the twenty-one compiled runs reach such a state (K, J,
E$^{\mathsf{B}}$, V, C), each first at a \textsc{Merge}.  All five
instantiate a polymorphic argument at the enclosing type variable
($f\,[X]$), so one cell's payload is another cell ($\beta \by \alpha$),
and a seal chain $\seal{X} ; \seal{Y}$ can outlive every frame that named
$X$ on a side.  No $\id{X}$ leaf failed in any run.

\end{document}
